\documentclass[../main.tex]{subfiles}
\begin{document}
\begin{longtable}{l p{14cm}}
	\hline
   \textbf{Thuật ngữ} & \textbf{Giải nghĩa} \\ \hline 
	\textbf{AsyncLocalStorage} & Cơ chế lưu trữ bối cảnh dữ liệu bất đồng bộ trong Node.js, giúp truyền tải thông tin định danh khách thuê tự động. \\
	\textbf{Deep Merge} & Thuật toán kết hợp đệ quy sâu dữ liệu để trộn lẫn bộ khung giao diện với cấu hình riêng của cửa hàng. \\
	\textbf{Master Template} & Bộ khung giao diện nền tảng tiêu chuẩn cung cấp bố cục hiển thị mặc định cho các cửa hàng. \\
	\textbf{Monorepo} & Kiến trúc quản lý mã nguồn lưu trữ và chia sẻ nhiều ứng dụng trong cùng một kho chứa duy nhất. \\
	\textbf{Multi-tenant} & Kiến trúc đa thuê bao cho phép một ứng dụng duy nhất phục vụ nhiều khách hàng khác nhau một cách độc lập. \\
	\textbf{NestJS} & Nền tảng phát triển ứng dụng máy chủ phụ trợ (backend) sử dụng TypeScript với kiến trúc mô-đun. \\
	\textbf{Next.js} & Nền tảng phát triển ứng dụng mặt tiền (frontend) dựa trên React, hỗ trợ kết xuất giao diện phía máy chủ. \\
	\textbf{Redis} & Hệ thống quản trị cơ sở dữ liệu lưu trữ trên bộ nhớ đệm giúp tối ưu hóa tốc độ phản hồi. \\
	\textbf{Tenant Config} & Dữ liệu cấu hình riêng biệt của từng khách thuê, ghi đè lên Master Template để tạo ra giao diện độc bản. \\
	\textbf{Tenant Owner} & Tác nhân chủ cửa hàng kinh doanh, người quản lý toàn bộ cấu hình, sản phẩm và giao dịch. \\
	\textbf{Turborepo} & Công cụ quản lý kho mã nguồn Monorepo, tối ưu hóa quá trình biên dịch thông qua lưu đệm tăng dần. \\
	\textbf{Webhook} & Cơ chế gọi ngược HTTP do hệ thống tự động kích hoạt để gửi thông báo thời gian thực khi có sự kiện thanh toán. \\
	\textbf{Zero-file} & Nguyên lý thiết kế giao diện trong đó cấu trúc hiển thị được định nghĩa bằng cấu hình JSON thay vì tệp mã nguồn. \\

    \hline
\end{longtable}

\end{document}